%
% Section 9.8_PowerSeries
%

%~~~~~~~~~~~~~~~~~~~~~~   []   ~~~~~~~~~~~~~~~~~~~~~~~~~~~
\begin{frame}
\frametitle{\Large {\color[rgb]{0.02,0.20,1.0} Power Series}}
\vspace{0.6cm}
\begin{block}{Homework Problems:}

\vspace{0.2cm}

TBA
%\hspace{0.2cm} {\bf 9.7}: 13, 15, 18, 19, 25, 26, 27 \hspace{0.5cm} {\bf 9.8}: 1-4, 11-22 \hspace{0.5cm} {\bf 9.10}: 21-26
	
\end{block}
\vspace{0.2cm}

\begin{block}{Section Example Problems to Read:}

\vspace{0.2cm}

\hspace{2.0cm} {\bf 9.7}: 3-6 \hspace{1.0cm} {\bf 9.8}: 1-7 \hspace{1.0cm} {\bf 9.10}: 3, 6
	
\end{block}
\vspace{0.1cm}

\begin{block}{This section will enable you to ...}
%-------LIST---------
\begin{itemize}
%	\item[...] define what power series is and where it is centered.
%	\item[...] explain how a function can be written as a power series.
	\item[...] determine the radius of convergence of a power series.
	\item[...] determine the interval of convergence of a power series.
	\item[...] find the Taylor polynomials of a given function.
	\item[...] find the Taylor polynomials from a known Taylor polynomial.
\end{itemize}
%-------------------
\end{block}
\vspace{3.0cm}
\end{frame}

%~~~~~~~~~~~~~~~~~~~~~~~~~~~~~~~~~~~~~~~~~~~~~~~~~~~~~~~~~
\begin{frame}
\frametitle{I. Introduction}
\scriptsize
\begin{columns}
\begin{column}{12.0cm}
Consider the following infinite series:
$$
\onslide<2->{f(\textcolor{blue}{x}) = } \sum \limits_{n=0}^{\infty} \onslide<1->{\textcolor<2->{blue}{x}}^n
$$
There is one major difference between this series and the ones we've been working with.  This series is a function of $x$ and for every value you plug into $f(x)$, you get an \onslide<3->{\textcolor{blue}{infinite series.}}

\begin{center}
	\setlength{\extrarowheight}{0.3cm}
	\begin{tabular}{| c | p{4.4cm} | c |}
		\hline
		$x=0$ & $\displaystyle{f(0) = \onslide<4->{\textcolor{blue}{\sum \limits_{n=1}^{\infty} (0)^n = \sum \limits_{n=1}^{\infty} 0 = 0}}}$ & \onslide<5->{\textcolor{blue}{converges (by the GST)}} \\
		\hline
		$x=1$ & $\displaystyle{f(1) = \onslide<5->{\textcolor{blue}{\sum \limits_{n=1}^{\infty} (1)^n}}}$ & \onslide<6->{\textcolor{blue}{diverges by the GST}}\\
		\hline
		$x=0.75$ & $\displaystyle{f(0.75) = \onslide<7->{\textcolor{blue}{\sum \limits_{n=1}^{\infty} (0.75)^n }}}$ & \onslide<8->{\textcolor{blue}{converges by the GST}}\\
		\hline
		$x=-\frac{1}{3}$ & $f\left(-\frac{1}{3}\right) = \onslide<9->{\textcolor{blue}{\displaystyle{\sum \limits_{n=1}^{\infty} \left(-\frac{1}{3}\right)^n}}}$ & \onslide<10->{\textcolor{blue}{converges by the GST}}\\
		\hline
		$x=10$ & $\displaystyle{f(10) = \onslide<11->{\textcolor{blue}{\sum \limits_{n=1}^{\infty} (10)^n}}}$ & \onslide<12->{\textcolor{blue}{diverges by the GST}}\\
		\hline
	\end{tabular}	
\end{center}

One obvious question then is: What values of $x$ make $f(x)$ converge?

\vspace{0.2cm}

\end{column}
\begin{column}{0.0cm}\end{column}
\end{columns}
\end{frame}

%~~~~~~~~~~~~~~~~~~~~~~~~~~~~~~~~~~~~~~~~~~~~~~~~~~~~~~~~~
\begin{frame}
\frametitle{II. Power Series {\normalsize({\em Section 9.8})}}
\scriptsize
\vspace{0.1cm}
\begin{columns}
\begin{column}{12.0cm}
	\begin{columns}\fcolorbox{blue}{white}{\begin{column}{10.4cm}
	% Type the contents of the box here
	\underline{Power Series}:
	A infinite series of the form
	$$
	f(x) = \sum \limits_{n=0}^{\infty} b_{n}(x-c)^n = b_0 + b_1(x-c) + b_2(x-c)^2 + ... + b_{n} (x-c)^n + ...
	$$
	is called a {\bf power series centered at $c$}, where $c$ is constant. \\
	\vspace{0.1cm}
	\end{column}}\begin{column}{0.0cm}\end{column}\end{columns}
	
	\vspace{0.5cm}
	
For each of the following power series, identify $b_{n}$ and $c$ in the power series definition. \\

\vspace{0.1cm}

\begin{columns}
\begin{column}{3.0cm}
	%--------EQN----------
	\begin{align*}
		g(x) = \sum \limits_{n=0}^{\infty} n^2 (x-3)^n \\
		\\
		\\
	\end{align*}
	%----------------------

	$b_{n} = \onslide<2->{\textcolor{blue}{n^2}}$ \quad and \quad $c = \onslide<3->{\textcolor{blue}{3}}$

	\vspace{1.2cm}
	
\end{column}
\begin{column}{4.0cm}
	%--------EQN----------
	\begin{align*}
	h(x) &= \sum \limits_{n=0}^{\infty} (-1)^n \frac{(x+1)^n}{n-1} \\
	&\onslide<5->{\textcolor{blue}{\,= \sum \limits_{n=0}^{\infty} \frac{(-1)^n}{n-1}\left(x-(-1)\right)^n}}
	\end{align*}
	%----------------------

	\vspace{0.2cm}
	
	$b_{n} = \onslide<6->{\textcolor{blue}{\displaystyle{\frac{(-1)^n}{n - 1}}}}$ \quad and \quad $c = \onslide<7->{\textcolor{blue}{-1}}$

	\vspace{1.2cm}
	
\end{column}
\begin{column}{3.5cm}
	%--------EQN----------
	\begin{align*}
		g(x) &= \sum \limits_{n=0}^{\infty} \frac{(2x)^n}{n^2} \\
		&\onslide<8->{\textcolor{blue}{\, = \sum \limits_{n=0}^{\infty} \frac{2^n x^n}{n^2}}} \onslide<9->{\textcolor{blue}{\, = \sum \limits_{n=0}^{\infty} \frac{2^n}{n^2}x^n}}
	\end{align*}
	%----------------------
	
	\vspace{0.2cm}
	
	$b_{n} = \onslide<10->{\textcolor{blue}{\displaystyle{\frac{2^n}{n^2}}}}$ \quad and \quad $c = \onslide<11->{\textcolor{blue}{0}}$

	\vspace{1.2cm}
	
\end{column}
\end{columns}

\end{column}
\begin{column}{0.0cm}\end{column}
\end{columns}
\end{frame}

%~~~~~~~~~~~~~~~~~~~~~~~~~~~~~~~~~~~~~~~~~~~~~~~~~~~~~~~~~
\begin{frame}
\frametitle{III. Interval and Radius of Convergence {\normalsize({\em Section 9.8})}}
\scriptsize
\vspace{0.1cm}
\begin{columns}
\begin{column}{12.0cm}
	\begin{columns}\fcolorbox{blue}{white}{\begin{column}{10.4cm}
	% Type the contents of the box here
	Consider the power series centered at $c$
	$$
	f(x) = \sum \limits_{n=0}^{\infty} b_{n}(x-c)^n = b_0 + b_1(x-c) + b_2(x-c)^2 + ... + b_{n} (x-c)^n + ...
	$$
	%\vspace{0.2cm}
	\underline{Interval of Convergence}:
	The interval that contains the values of $x$ for which $f(x)$ is a convergent infinite series is called the {\bf interval of convergence}.  \\
	\vspace{0.2cm}
	\underline{Radius of Convergence}: 
	Half the width of the interval of convergence.\\
	\end{column}}\begin{column}{0.0cm}\end{column}\end{columns}
	\vspace{0.3cm}
	\centering
	There are 3 main forms of the interval of convergence.  They are
	\vspace{0.2cm}
	\setlength{\extrarowheight}{0.07cm}
	\begin{tabular}{| c | c | c |}
		\hline
		Interval of Convergence & Number Line & Radius of Convergence \\
		\hline
		\multirow{3}{*}{$c$} & \multirow{3}{*}{\includegraphics[scale = 0.125]{Figures/9_8_Fig1.png}} & \multirow{3}{*}{0}\\
		&&\\
		&&\\
		\hline
		\multirow{3}{*}{$(c-R,c+R)$} & \multirow{3}{*}{\includegraphics[scale = 0.125]{Figures/9_8_Fig2.png}} & \multirow{3}{*}{R}\\
		&&\\
		&&\\
		\hline
		\multirow{3}{*}{$(-\infty, \infty)$} & \multirow{3}{*}{\includegraphics[scale = 0.125]{Figures/9_8_Fig3.png}} & \multirow{3}{*}{$\infty$}\\
		&&\\
		&&\\
		\hline
	\end{tabular}
\end{column}
\begin{column}{0.0cm}\end{column}
\end{columns}
\end{frame}

%~~~~~~~~~~~~~~~~~~~~~~~~~~~~~~~~~~~~~~~~~~~~~~~~~~~~~~~~~
\begin{frame}
\frametitle{IV. Finding the Interval of Convergence {\normalsize({\em Section 9.8})}}
\scriptsize
\begin{columns}
\begin{column}{12.0cm}
	To find the interval of convergence of a power series
	$$
	f(x) = \sum \limits_{n=0}^{\infty} b_{n}(x-c)^n
	$$
	we go through the following steps:\\
	\vspace{0.2cm}
	%-------LIST---------
	\begin{itemize}
		\item[1.] Identify \onslide<2->{\textcolor{blue}{$c$ and $b_{n}$.}}
		\item[2.] Find the radius of convergence, $R$, using \onslide<3->{\textcolor{blue}{$\displaystyle{\frac{1}{R} = \lim \limits_{n \to \infty}\, \left|\frac{b_{n+1}}{b_{n}}\right|}$ \quad (ratio test limit)}}
		\item[] $\circ$ If \,\,\onslide<4->{\textcolor{blue}{$R = 0$}}\,\, then the interval of convergence is \onslide<5->{\textcolor{blue}{the single point $c$.}}
		\item[] $\circ$ If \,\,\onslide<6->{\textcolor{blue}{$R = \infty$}} \,\,then the interval of convergence is \,\,\onslide<7->{\textcolor{blue}{$(-\infty,\infty)$.}}
		\item[] $\circ$ If \,\,\onslide<8->{\textcolor{blue}{$0 < R < \infty$ (finite)}}\,\, then the interval of convergence is \onslide<9->{\textcolor{blue}{at least \,\,$(c-R,c+R)$.}} \\
		\vspace{0.2cm}
		%-------LIST---------
		\begin{itemize}
		\scriptsize
			\item if $f(c-R)$ converges then we \onslide<10->{\textcolor{blue}{include the left endpoint of $(c-R,c+R)$}}
			\vspace{0.1cm}
			\begin{center}
				\onslide<11->{\textcolor{blue}{square bracket $\displaystyle{\rightarrow [c-R,c+R)}$}}
			\end{center}
			\vspace{0.1cm}
			\item if $f(c+R)$ converges then we \onslide<12->{\textcolor{blue}{include the right endpoint of $(c-R,c+R)$}}
			\vspace{0.1cm}
			\begin{center}
				\onslide<13->{\textcolor{blue}{$\displaystyle{(c-R,c+R] \leftarrow }$ square bracket}}
			\end{center}
			\vspace{0.2cm}
				\item if both series above converge then we \onslide<14->{\textcolor{blue}{include both endpoints}}
				\vspace{0.2cm}
			\begin{center}
				\onslide<15->{\textcolor{blue}{square bracket $\displaystyle{\rightarrow [c-R,c+R] \leftarrow }$ square bracket}}
			\end{center}
		\end{itemize}
		%------------------- 
	\end{itemize}
	%-------------------
\end{column}
\begin{column}{0.0cm}\end{column}
\end{columns}
\end{frame}

%~~~~~~~~~~~~~~~~~~~~~~~~~~~~~~~~~~~~~~~~~~~~~~~~~~~~~~~~~
\begin{frame}
\frametitle{}
\scriptsize

\vspace{0.2cm}

\begin{columns}
\begin{column}{12.0cm}
Find the interval of convergence of the power series \\

\begin{columns}
\begin{column}{4.0cm}
	$\displaystyle{\sum \limits_{n=0}^{\infty} \left(\frac{x}{5}\right)^n \onslide<2->{\textcolor{blue}{\,= \sum \limits_{n=0}^{\infty} \frac{x^n}{5^n}}}}$
\end{column}
\begin{column}{6.5cm}
	\underline{Step 1}: Identify \quad $c = \onslide<3->{\textcolor{blue}{0}}$ \quad and \quad $\displaystyle{b_{n} = \onslide<4->{\textcolor{blue}{\frac{1}{5^n}}}}$ \\
	\underline{Step 2}: Compute \onslide<1->{\textcolor{blue}{$\displaystyle{\frac{1}{R} = \lim \limits_{n \to \infty}\, \left|\frac{b_{n+1}}{b_{n}}\right|}$}}\\
\end{column}
\end{columns}

%--------EQN----------
\begin{align*}
	\frac{1}{R} = \lim \limits_{n \to \infty}\, \left|\frac{b_{n+1}}{b_{n}} \right|	= 
	\onslide<6->{\textcolor{blue}{\lim \limits_{n \to \infty}\, \left|\frac{1}{5^{n+1}}\cdot\frac{5^n}{1}\right|}} 
	\onslide<7->{\textcolor{blue}{\,= \lim \limits_{n \to \infty}\, \left|\frac{1}{\only<2->{\cancel}{5^{n}}\cdot5}\cdot\frac{\only<2->{\cancel}{5^{n}}}{1}\right|}}
	\onslide<8->{\textcolor{blue}{\,= \lim \limits_{n \to \infty}\, \frac{1}{5}}}
	\onslide<9->{\textcolor{blue}{\,= \frac{1}{5}}} \hspace{2.0cm}
\end{align*}
%----------------------
So \quad \onslide<10->{\textcolor{blue}{$\displaystyle{\frac{1}{R} = \frac{1}{5}}$}} \quad and the radius of convergence (RoC) is \quad \onslide<11->{\textcolor{blue}{$\displaystyle{R = 5}$.}}  \\

\vspace{0.2cm}

Therefore, the interval of convergence (IoC) is at least \onslide<12->{\textcolor{blue}{$(c-R,c+R) = (-5,5)$}}

\vspace{0.2cm}

Check Left Endpoint: \onslide<13->{\textcolor{blue}{$x = -5$}} \\

\vspace{0.2cm}

$\displaystyle{\onslide<14->{\textcolor{blue}{\sum \limits_{n=0}^{\infty} \left(\frac{(-5)}{5}\right)^n}} \onslide<15->{\textcolor{blue}{\,= \sum \limits_{n=0}^{\infty} \left(-1\right)^n}}}$ \quad 
\onslide<16->{\textcolor{blue}{diverges by the GST $\rightarrow |r| = 1 \ge 1$ }}

\vspace{0.3cm}

Check Right Endpoint: \onslide<17->{\textcolor{blue}{$x = 5$}} \\

\vspace{0.2cm}

$\displaystyle{\onslide<18->{\textcolor{blue}{\sum \limits_{n=0}^{\infty} \left(\frac{(5)}{5}\right)^n}} \onslide<19->{\textcolor{blue}{\,= \sum \limits_{n=0}^{\infty} \left(1\right)^n}}}$ \quad 
\onslide<20->{\textcolor{blue}{diverges by the GST $\rightarrow |r| = 1 \ge 1$ }}

\vspace{0.3cm}

Therefore, the \fbox{IoC = \onslide<21->{\textcolor{blue}{$(-5,5)$}}}
\vspace{10.5cm}

\end{column}
\begin{column}{0.0cm}
	
\end{column}
\end{columns}
\end{frame}

%~~~~~~~~~~~~~~~~~~~~~~~~~~~~~~~~~~~~~~~~~~~~~~~~~~~~~~~~~
\begin{frame}
\frametitle{}
\scriptsize

\vspace{0.2cm}

\begin{columns}
\begin{column}{12.0cm}

\begin{columns}
\begin{column}{4.8cm}
	$\displaystyle{\sum \limits_{n=0}^{\infty} \frac{(-1)^n(x+1)^n}{n} \onslide<2->{\textcolor{blue}{\,= \sum \limits_{n=0}^{\infty} \frac{(-1)^n}{n}(x+1)^n}}}$\\
\end{column}
\begin{column}{6.0cm}
	\underline{Step 1}: Identify \quad $c = \onslide<3->{\textcolor{blue}{-1}}$ \quad and \quad $\displaystyle{b_{n} = \onslide<4->{\textcolor{blue}{\frac{(-1)^n}{n}}}}$ \\
	\underline{Step 2}: Compute \onslide<1->{\textcolor{blue}{$\displaystyle{\frac{1}{R} = \lim \limits_{n \to \infty}\, \left|\frac{b_{n+1}}{b_{n}}\right|}$}}\\
\end{column}
\end{columns}

%--------EQN----------
\begin{align*}
	\frac{1}{R} = \lim \limits_{n \to \infty}\, \left|\frac{b_{n+1}}{b_{n}} \right|	= 
	\onslide<6->{\textcolor{blue}{\lim \limits_{n \to \infty}\, \left|\frac{(-1)^{n+1}}{n+1}\cdot\frac{n}{(-1)^n}\right|}}
	\onslide<7->{\textcolor{blue}{\, = \lim \limits_{n \to \infty}\, \frac{n}{n+1}}}
	\onslide<8->{\textcolor{blue}{\, \overset{LH}{=} \lim \limits_{n \to \infty}\, \frac{1}{1}}} 
	\onslide<9->{\textcolor{blue}{\,= 1}} \hspace{2.0cm}
\end{align*}
%----------------------
So \quad \onslide<10->{\textcolor{blue}{$\displaystyle{\frac{1}{R} = 1}$}} \quad and the radius of convergence (RoC) is \quad \onslide<11->{\textcolor{blue}{$\displaystyle{R = 1}$.}}  \\

\vspace{0.2cm}

Therefore, the interval of convergence (IoC) is at least \onslide<12->{\textcolor{blue}{$(c-R,c+R) = (-2,0)$}}

\vspace{0.2cm}

Check Left Endpoint: \onslide<13->{\textcolor{blue}{$x = -2$}} \\

\vspace{0.2cm}

$\displaystyle{\onslide<14->{\textcolor{blue}{\sum \limits_{n=0}^{\infty} \frac{(-1)^n((-2)+1)^n}{n}}} \onslide<15->{\textcolor{blue}{\,= \sum \limits_{n=0}^{\infty} \frac{(-1)^n(-1)^n}{n}}} \onslide<16->{\textcolor{blue}{\,= \sum \limits_{n=0}^{\infty} \frac{((-1)^n)^2}{n}}} \onslide<17->{\textcolor{blue}{\,= \sum \limits_{n=0}^{\infty} \frac{1}{n}}} }$ \quad 
\onslide<18->{\textcolor{blue}{diverges by the p-test }}

\vspace{0.3cm}

Check Right Endpoint: \onslide<19->{\textcolor{blue}{$x = 0$}} \\

\vspace{0.2cm}

$\displaystyle{\onslide<20->{\textcolor{blue}{\sum \limits_{n=0}^{\infty} \frac{(-1)^n((0)+1)^n}{n}}} \onslide<21->{\textcolor{blue}{\,= \sum \limits_{n=0}^{\infty} \frac{(-1)^n(1)^n}{n}}} \onslide<22->{\textcolor{blue}{\,= \sum \limits_{n=0}^{\infty} \frac{(-1)^n}{n}}} }$ \quad 
\onslide<23->{\textcolor{blue}{ Apply AST $\implies$ let  $a_n = \frac{1}{n}$. }}\\
\onslide<24->{\textcolor{blue}{Since \, (i) \, $\displaystyle{\frac{1}{n}}$ is clearly decreasing and}} \onslide<25->{\textcolor{blue}{(ii) \,$\displaystyle{\lim \limits_{n \to \infty}\, \frac{1}{n} = 0}$}} \onslide<26->{\textcolor{blue}{, \, then $\displaystyle{\sum \limits_{n=0}^{\infty} \frac{(-1)^n}{n}}$ converges by AST }}

\vspace{0.1cm}

Therefore, the \fbox{IoC = \onslide<27->{\textcolor{blue}{$(-2,0]$}}}
\vspace{10.5cm}

\end{column}
\begin{column}{0.0cm}
	
\end{column}
\end{columns}
\end{frame}

%~~~~~~~~~~~~~~~~~~~~~~~~~~~~~~~~~~~~~~~~~~~~~~~~~~~~~~~~~
\begin{frame}
\frametitle{}
\scriptsize

\vspace{0.0cm}

\begin{columns}
\begin{column}{12.5cm}

\begin{columns}
\begin{column}{4.0cm}
	$\displaystyle{\sum \limits_{n=0}^{\infty} \frac{n!(x-3)^n}{3^n}}$
\end{column}
\begin{column}{6.0cm}
	\underline{Step 1}: Identify \quad $c = \onslide<2->{\textcolor{blue}{3}}$ \quad and \quad $\displaystyle{b_{n} = \onslide<3->{\textcolor{blue}{\frac{n!}{3^n}}}}$ \\
	\underline{Step 2}: Compute \onslide<1->{\textcolor{blue}{$\displaystyle{\frac{1}{R} = \lim \limits_{n \to \infty}\, \left|\frac{b_{n+1}}{b_{n}}\right|}$}}
\end{column}
\end{columns}

\vspace{0.4cm}

%--------EQN----------
\begin{align*}
	\frac{1}{R} = \lim \limits_{n \to \infty}\, \left|\frac{b_{n+1}}{b_{n}} \right|	= 
	\onslide<5->{\textcolor{blue}{\lim \limits_{n \to \infty}\, \left|\frac{(n+1)!}{3^{n+1}}\cdot\frac{3^n}{n!}\right|}} 
	\onslide<6->{\textcolor{blue}{\,= \lim \limits_{n \to \infty}\, \left|\frac{(n+1)\only<7->{\cancel}{n!}}{\only<7->{\cancel}{3^{n}}\cdot3^{1}}\cdot\frac{\only<7->{\cancel}{3^{n}}}{\only<7->{\cancel}{n!}}\right|}} 
	\onslide<8->{\textcolor{blue}{\,= \lim \limits_{n \to \infty}\, \frac{n+1}{3}}}
	\onslide<9->{\textcolor{blue}{\,= \infty}}
\end{align*}
%----------------------

\vspace{0.2cm}

So \quad \onslide<10->{\textcolor{blue}{$\displaystyle{\frac{1}{R} = \infty}$}} \quad and the radius of convergence (RoC) is \quad \onslide<11->{\textcolor{blue}{$\displaystyle{R = 0}$.}}

\vspace{0.4cm}

Therefore, \fbox{IoC = \onslide<12->{\textcolor{blue}{$c = 3$}}}

\vspace{5.2cm}

\end{column}
\begin{column}{0.0cm}
	
\end{column}
\end{columns}
\end{frame}

%~~~~~~~~~~~~~~~~~~~~~~~~~~~~~~~~~~~~~~~~~~~~~~~~~~~~~~~~~
\begin{frame}
\frametitle{}
\scriptsize

\vspace{0.2cm}

\begin{columns}
\begin{column}{12.0cm}

\begin{columns}
\begin{column}{4.0cm}
	$\displaystyle{\sum \limits_{n=0}^{\infty} \frac{x^n}{(2n)!}}$
\end{column}
\begin{column}{6.5cm}
	\underline{Step 1}: Identify \quad $c = \onslide<2->{\textcolor{blue}{0}}$ \quad and \quad $\displaystyle{b_{n} = \onslide<3->{\textcolor{blue}{\frac{1}{(2n)!}}}}$ \\
	\underline{Step 2}: Compute \onslide<4->{\textcolor{blue}{$\displaystyle{\frac{1}{R} = \lim \limits_{n \to \infty}\, \left|\frac{b_{n+1}}{b_{n}}\right|}$}}\\
\end{column}
\end{columns}

\vspace{0.2cm}

%--------EQN----------
\begin{align*}
	\frac{1}{R} = \lim \limits_{n \to \infty}\, \left|\frac{b_{n+1}}{b_{n}} \right|	&\,=
	\onslide<5->{\textcolor{blue}{\lim \limits_{n \to \infty}\, \left|\frac{1}{(2(n+1))!}\cdot\frac{(2n)!}{1}\right|}}
	\onslide<6->{\textcolor{blue}{\, =\lim \limits_{n \to \infty}\, \left|\frac{1}{(2n+2)!}\cdot\frac{(2n)!}{1}\right|}} \\
	&\,\onslide<7->{\textcolor{blue}{\, =\lim \limits_{n \to \infty}\, \left|\frac{1}{(2n+2)(2n+1)\only<8->{\cancel}{(2n)!}}\cdot\frac{\only<8->{\cancel}{(2n)!}}{1}\right|}}
	\onslide<9->{\textcolor{blue}{\, =\lim \limits_{n \to \infty}\,\frac{1}{(2n+2)(2n+1)}}} 
	\onslide<10->{\textcolor{blue}{\,= 0}}
\end{align*}
%----------------------

\vspace{0.2cm}

So \quad \onslide<11->{\textcolor{blue}{$\displaystyle{\frac{1}{R} = 0}$}} \quad and the radius of convergence (RoC) is \quad \onslide<12->{\textcolor{blue}{$\displaystyle{R = \infty}$.}}  \\

\vspace{0.4cm}

Therefore, \fbox{IoC = \onslide<13->{\textcolor{blue}{$(-\infty,\infty)$}}}

\vspace{5.5cm}

\end{column}
\begin{column}{0.0cm}
	
\end{column}
\end{columns}
\end{frame}


%~~~~~~~~~~~~~~~~~~~~~~~~~~~~~~~~~~~~~~~~~~~~~~~~~~~~~~~~~
\begin{frame}
\frametitle{V. Taylor Series \& Maclaurin Series {\normalsize({\em Section 9.7})}}
\scriptsize
\begin{columns}
\begin{column}{12.0cm}
\vspace{0.2cm}

There is a very important power series known as a Taylor Series.

\vspace{0.2cm}

\begin{columns}\fcolorbox{blue}{white}{\begin{column}{10.4cm}
% Type the contents of the box here
The {\bf Taylor series for ${\bf f(x)}$ at c} is defined as
$$
\sum \limits_{n=0}^{\infty} \frac{f^{(n)}(c)}{n!}(x-c)^n = f(c) + f'(c)(x-c) + \frac{f''(c)}{2}(x-c)^2 + \frac{f^{(3)}(c)}{3!}(x-c)^3 + ...
$$
If $c=0$, then we call this series a {\bf Maclaurin series for ${\bf f(x)}$}
\end{column}}\begin{column}{0.0cm}\end{column}\end{columns}

\vspace{0.8cm}

When discussing finite sum versions of the above, we defined the following

\vspace{0.2cm}

\begin{columns}\fcolorbox{blue}{white}{\begin{column}{10.4cm}
% Type the contents of the box here
The {\bf Taylor series for ${\bf f(x)}$ at c of degree K} is defined as
$$
\sum \limits_{n=0}^{K} \frac{\only<2->{\textcolor{blue}}{f^{(n)}(c)}}{n!}(x-c)^n = \only<2->{\textcolor{blue}}{f(c)} + \only<2->{\textcolor{blue}}{f'(c)}(x-c) + \frac{\only<2->{\textcolor{blue}}{f''(c)}}{2}(x-c)^2 + ... + \frac{\only<2->{\textcolor{blue}}{f^{(K)}(c)}}{K!}(x-c)^K
$$
If $c=0$, then we call this series a {\bf Maclaurin series for ${\bf f(x)}$ of degree K}
\end{column}}\begin{column}{0.0cm}\end{column}\end{columns}

\vspace{0.2cm}

To compute Taylor series we need to find:\quad \only<2->{\textcolor{blue}}{$f(c),\, f'(c), \, f''(c),\, ...,\, f^{(K)}(c)$}\quad and plug these into the above formula.
\vspace{2.2cm}

\end{column}
\begin{column}{0.0cm}\end{column}
\end{columns}
\end{frame}

%~~~~~~~~~~~~~~~~~~~~~~~~~~~~~~~~~~~~~~~~~~~~~~~~~~~~~~~~~
\begin{frame}
\frametitle{}
\scriptsize

\vspace{0.2cm}

\begin{columns}
\begin{column}{12.0cm}
For \quad $\displaystyle{f(x) = e^{-2x}}$ \quad find \\

(i) the Taylor series of degree 3 centered at c = 1; \\
\vspace{0.02cm}
(ii) the Maclaurin series of degree 3. \\

\vspace{0.1cm}

Strategy for (i):\\
\vspace{0.05cm}
%-------LIST---------
\begin{itemize}
	\item compute \, \onslide<2->{\textcolor{blue}{$f(1), \, f'(1), \, f''(1), \, f'''(1)$}}
	\item plug into \onslide<3->{\textcolor{blue}{\quad $\displaystyle{f(1) + f'(1)(x-1) + \frac{f''(1)}{2}(x-1)^2 + \frac{f'''(1)}{6}(x-1)^3}$}}
\end{itemize}
%-------------------
	\setlength{\extrarowheight}{0.2cm}
	\begin{tabular}{p{0.3cm}lclcl}
		&$\displaystyle{f(x) = 
		\onslide<4->{\textcolor{blue}{e^{-2x}}}}$ & 
		\onslide<5->{\textcolor{blue}{$\implies$}} & 
		\onslide<5->{\textcolor{blue}{$\displaystyle{f(1) = e^{-2}}$}} & 
		\onslide<15->{\textcolor{blue}{\&}} & 
		\onslide<15->{\textcolor{blue}{$\displaystyle{f(0) = 1}$}} \\
		&$\displaystyle{f'(x) = 
		\onslide<6->{\textcolor{blue}{-2e^{-2x}}}}$ & 
		\onslide<7->{\textcolor{blue}{$\implies$}} & 
		\onslide<7->{\textcolor{blue}{$\displaystyle{f'(1) = -2e^{-2}}$}} & 
		\onslide<16->{\textcolor{blue}{\&}} & 
		\onslide<16->{\textcolor{blue}{$\displaystyle{f'(0) = -2}$}} \\
		&$\displaystyle{f''(x) = 
		\onslide<8->{\textcolor{blue}{-2(-2)e^{-2x} = 4e^{-2x}}}}$ & 
		\onslide<9->{\textcolor{blue}{$\implies$}} & 
		\onslide<9->{\textcolor{blue}{$\displaystyle{f''(1) = 4e^{-2}}$}} & 
		\onslide<17->{\textcolor{blue}{\&}} & 
		\onslide<17->{\textcolor{blue}{$\displaystyle{f''(0) = 4}$}} \\
		&$\displaystyle{f'''(x) = 
		\onslide<10->{\textcolor{blue}{4(-2)e^{-2x} = -8e^{-2x}}}}$ & 
		\onslide<11->{\textcolor{blue}{$\implies$}} & 
		\onslide<11->{\textcolor{blue}{$\displaystyle{f'''(1) = -8e^{-2}}$}} & 
		\onslide<18->{\textcolor{blue}{\&}} & 
		\onslide<18->{\textcolor{blue}{$\displaystyle{f''(0) = -8}$}}
	\end{tabular}

\vspace{0.1cm}

So the 3rd degree Taylor series is \onslide<12->{\textcolor{blue}{\fbox {\quad $\displaystyle{e^{-2} - 2e^{-2}(x-1) + \frac{4e^{-2}}{2}(x-1)^2 - \frac{8e^{-2}}{6}(x-1)^3}$ \quad }}} \\

\vspace{0.15cm}
Strategy for (ii):\\
\vspace{0.05cm} 
%-------LIST---------
\begin{itemize}
	\item compute \, \onslide<13->{\textcolor{blue}{$f(0), \, f'(0), \, f''(0), \, f'''(0)$}}
	\item plug into \onslide<14->{\textcolor{blue}{\quad $\displaystyle{f(0) + f'(0)x + \frac{f''(0)}{2}x^2 + \frac{f'''(0)}{6}x^3}$}}
\end{itemize}
%-------------------

So the 3rd degree Maclaurin series is \onslide<19->{\textcolor{blue}{\fbox{\quad $\displaystyle{1 - 2x + \frac{4}{2}x^2 - \frac{8}{6}x^3}$ \quad}}} \\

\vspace{3.2cm}

\end{column}
\begin{column}{0.0cm}
	
\end{column}
\end{columns}
\end{frame}

%~~~~~~~~~~~~~~~~~~~~~~~~~~~~~~~~~~~~~~~~~~~~~~~~~~~~~~~~~
\begin{frame}
\frametitle{}
\scriptsize

\vspace{0.2cm}

\begin{columns}
\begin{column}{12.0cm}
For \quad $\displaystyle{f(x) = \sin x}$ \quad find the Taylor series of degree 2 centered at $c = \frac{\pi}{4}$; \\

\vspace{0.1cm}

Strategy:\\
\vspace{0.05cm}
%-------LIST---------
\begin{itemize}
	\item compute \, \onslide<2->{\textcolor{blue}{$f\left(\frac{\pi}{4}\right), \, f'\left(\frac{\pi}{4}\right), \, f''\left(\frac{\pi}{4}\right)$}}
	\item plug into \onslide<3->{\textcolor{blue}{\quad $\displaystyle{f\left(\frac{\pi}{4}\right) + f'\left(\frac{\pi}{4}\right)\left(x-\frac{\pi}{4}\right) + \frac{f''\left(\frac{\pi}{4}\right)}{2}\left(x-\frac{\pi}{4}\right)^2}$}}
\end{itemize}
%-------------------
	\setlength{\extrarowheight}{0.4cm}
	\begin{tabular}{p{0.3cm}lcl}
		&$\displaystyle{f(x) = \onslide<4->{\textcolor{blue}{\sin x}}}$ & 
		\onslide<5->{\textcolor{blue}{$\implies$}} & 
		\onslide<5->{\textcolor{blue}{$\displaystyle{f\left(\frac{\pi}{4}\right) = \sin \left(\frac{\pi}{4}\right) = \frac{\sqrt{2}}{2}}$}}\\
		&$\displaystyle{f'(x) = 
		\onslide<6->{\textcolor{blue}{\cos x}}}$ & 
		\onslide<7->{\textcolor{blue}{$\implies$}} & 
		\onslide<7->{\textcolor{blue}{$\displaystyle{f'\left(\frac{\pi}{4}\right) = \cos\left(\frac{\pi}{4}\right) = \frac{\sqrt{2}}{2}}$}} \\
		&$\displaystyle{f''(x) = 
		\onslide<8->{\textcolor{blue}{-\sin x}}}$ & 
		\onslide<9->{\textcolor{blue}{$\implies$}} & 
		\onslide<9->{\textcolor{blue}{$\displaystyle{f''\left(\frac{\pi}{4}\right) = - \frac{\sqrt{2}}{2}}$}}
	\end{tabular}

\vspace{0.5cm}

So the 2nd degree Taylor series at $\frac{\pi}{4}$ is \quad \onslide<10->{\textcolor{blue}{\fbox {\quad $\displaystyle{\frac{\sqrt{2}}{2} + \frac{\sqrt{2}}{2}\left(x-\frac{\pi}{4}\right) - \frac{\sqrt{2}}{4}\left(x-\frac{\pi}{4}\right)^2}$ \quad }}} \\

\vspace{3.2cm}

\end{column}
\begin{column}{0.0cm}
	
\end{column}
\end{columns}
\end{frame}

%~~~~~~~~~~~~~~~~~~~~~~~~~~~~~~~~~~~~~~~~~~~~~~~~~~~~~~~~~
\begin{frame}
\frametitle{VI. Finding Taylor series from other series {\normalsize({\em Section 9.10})}}
\scriptsize
\begin{columns}
\begin{column}{12.0cm}
	On your formula sheets you will find a list of Maclaurin Series such as:
	%--------EQN----------
	\begin{align*}
		e^x &\,= 1 + x + \frac{x^2}{2!} + \frac{x^3}{3!} + ... \\
		\sin x &\,= x - \frac{x^3}{3!} + \frac{x^5}{5!} - ... \\
		\cos x &\,= 1 - \frac{x^2}{2!} + \frac{x^4}{4!} - ... 
	\end{align*}
	%----------------------
	
	\vspace{0.4cm}
	
	We can use these to find the Maclaurin series of other functions such as:
    %--------EQN----------
	\begin{align*}
		e^{x^2} &\,= 1 + (\,\onslide<2->{\textcolor{blue}{x^2}}\,) + \frac{(\,\onslide<2->{\textcolor{blue}{x^2}}\,)^2}{2!} + \frac{(\,\onslide<2->{\textcolor{blue}{x^2}}\,)^3}{3!} + ... = \onslide<3->{\textcolor{blue}{ 1 + x^2 + \frac{x^4}{2!} + \frac{x^6}{3!} + ... }}\\
		\sin (-2x) &\,= (\,\onslide<4->{\textcolor{blue}{-2x}}\,) - \frac{(\,\onslide<4->{\textcolor{blue}{-2x}}\,)^3}{3!} + \frac{(\,\onslide<4->{\textcolor{blue}{-2x}}\,)^5}{5!} - ... = \onslide<5->{\textcolor{blue}{ -2x + \frac{8x^3}{3!} - \frac{32x^5}{5!} - ...}} \\
		\cos (\sqrt{x}) &\,= 1 - \frac{(\,\onslide<6->{\textcolor{blue}{\sqrt{x}}}\,)^2}{2!} + \frac{(\,\onslide<6->{\textcolor{blue}{\sqrt{x}}}\,)^4}{4!} - ... = \onslide<7->{\textcolor{blue}{ 1 - \frac{x}{2!} + \frac{x^2}{4!} - ...}} \qquad (x\ge0)\\
	\end{align*}
	%----------------------
\end{column}
\begin{column}{0.0cm}\end{column}
\end{columns}
\end{frame}
